\documentclass[border=10pt]{standalone}
\usepackage{fontspec}
\setmainfont{Noto Serif}[Ligatures=TeXOff]
\XeTeXgenerateactualtext=1
\newfontfamily\bengalifont{Noto Serif Bengali}[Script=Bengali, NFSSFamily=bnfam]
\usepackage[Latin,Bengali]{ucharclasses}
% group-free font switching (as in header-common.tex): robust inside braces, cells and quotes
\makeatletter
\def\BN@saved{\rmdefault}
\def\BN@bn{bnfam}
\newcommand\BNin{\edef\BN@tmp{\f@family}\ifx\BN@tmp\BN@bn\else\let\BN@saved\BN@tmp\fi\fontfamily{bnfam}\selectfont}
\newcommand\BNout{\fontfamily{\BN@saved}\selectfont}
\makeatother
\setTransitionsFor{Bengali}{\BNin}{\BNout}
% Typographic quotes and dashes are not in a ucharclasses class (class 0), so a Bangla word right
% after “ or ‘ would not trigger the Bangla font. Treat them as Latin so that every switch fires.
\XeTeXcharclass"2018=\XeTeXcharclass`A
\XeTeXcharclass"2019=\XeTeXcharclass`A
\XeTeXcharclass"201C=\XeTeXcharclass`A
\XeTeXcharclass"201D=\XeTeXcharclass`A
\XeTeXcharclass"2013=\XeTeXcharclass`A
\XeTeXcharclass"2014=\XeTeXcharclass`A
\XeTeXcharclass"2026=\XeTeXcharclass`A
\usepackage{newunicodechar}
\newunicodechar{।}{{\bengalifont\char"0964}}
\newunicodechar{॥}{{\bengalifont\char"0965}}

\usepackage[table]{xcolor}
\usepackage{booktabs,tabularx,array}
\definecolor{bapilrule}{HTML}{1F2937}
\definecolor{bapilstripe}{HTML}{F3F4F6}
\renewcommand{\arraystretch}{1.28}
\setlength{\tabcolsep}{5pt}

\begin{document}
\begin{minipage}{17cm}
{\large\bfseries সারণি ১৪।}\ {\large শাসন-কাঠামোর সংস্থাসমূহ, গঠন ও সিদ্ধান্তের অধিকার}\par\vspace{7pt}
\small
\rowcolors{2}{bapilstripe}{white}
\begin{tabularx}{\linewidth}{@{}>{\hsize=0.782\hsize\raggedright\arraybackslash}X>{\hsize=1.195\hsize\raggedright\arraybackslash}X>{\hsize=1.287\hsize\raggedright\arraybackslash}X>{\hsize=0.736\hsize\raggedright\arraybackslash}X@{}}
\arrayrulecolor{bapilrule}\toprule
\rowcolor{white}\textbf{সংস্থা} & \textbf{গঠন} & \textbf{প্রধান দায়িত্ব} & \textbf{সিদ্ধান্তের অধিকার} \\
\midrule
\textbf{স্টিয়ারিং বোর্ড} & ৭ থেকে ৯ জন সদস্য: আইসিটি বিভাগ (সহ-সভাপতি), স্বাগতিক প্রতিষ্ঠান (সহ-সভাপতি), পরিকল্পনা কমিশন, একটি ডেটা-রক্ষক সংস্থা, একাডেমিয়া, নাগরিক সমাজ, শিল্পখাত; BAPIL পরিচালক সদস্য-সচিব & কৌশল, বার্ষিক পরিকল্পনা ও বাজেট; অংশীদারিত্ব অনুমোদন; স্বাগতিক সংস্থার সঙ্গে যৌথভাবে G4 গ্রহণ-সিদ্ধান্ত & পরিকল্পনা ও বাজেট অনুমোদন \\
\textbf{নৈতিকতা ও ঝুঁকি পর্যালোচনা কমিটি (ERRC)} & ৫ থেকে ৭ জন, সংখ্যাগরিষ্ঠ বহিঃস্থ: আইন ও মানবাধিকার, ডেটা সুরক্ষা, নৈতিকতা, বিষয়-বিশেষজ্ঞ, নাগরিক সমাজ, কারিগরি পর্যালোচক & G1, G2 ও G4-এ T2 ও T3 স্তরের ব্যবস্থা পর্যালোচনা; ঘটনা পর্যালোচনা; সিদ্ধান্তের সারাংশ প্রকাশ & পাইলট ও প্রয়োগে ভেটো, যা BAPIL-এর জন্য এবং স্বাগতিকের সংবিধির মধ্যে স্বাগতিকের জন্য বাধ্যতামূলক (অনুচ্ছেদ ৫.৩) \\
\textbf{কারিগরি স্থাপত্য বোর্ড (TAB)} & প্রধান প্রকৌশলী, নিরাপত্তা প্রধান, দুইজন বহিঃস্থ কারিগরি পর্যালোচক & স্থাপত্য ও নিরাপত্তা মানদণ্ড; নকশা পর্যালোচনা; মডেল ও নির্ভরতা অনুমোদন & G1-এ নকশা অনুমোদন \\
\textbf{আন্তর্জাতিক উপদেষ্টা প্যানেল} & ব্যক্তিগত সক্ষমতায় আমন্ত্রিত বিশেষজ্ঞ & পদ্ধতি, বেঞ্চমার্ক ও শাসনব্যবস্থার সমকক্ষ পর্যালোচনা; আন্তর্জাতিক সংযোগ & শুধু পরামর্শমূলক \\
\textbf{পরিচালক ও ইউনিটসমূহ} & পরিচালক, হোয়াইট পেপারের ছয়টি বিষয়ভিত্তিক ইউনিট এবং একটি আন্তঃবিভাগীয় প্ল্যাটফর্ম প্রকৌশল দল & গবেষণা, পাইলট, প্রশিক্ষণ ও প্রকাশনা বাস্তবায়ন & অনুমোদিত পরিকল্পনার মধ্যে পরিচালনাগত সিদ্ধান্ত \\
\textbf{স্বাগতিক সংস্থার প্রোডাক্ট ওনার} & পাইলট পরিচালনাকারী প্রতিটি সংস্থার মনোনীত কর্মকর্তা & চাহিদা নির্ধারণ, পাইলট ব্যবহারকারী সংগ্রহ, G3 ও G4-এ গ্রহণ বা প্রত্যাখ্যান & নিজ সংস্থায় পাইলট গ্রহণ \\
\bottomrule
\end{tabularx}
\par\vspace{3pt}{\scriptsize\color{gray} উৎস: BAPIL টেকনিক্যাল ব্লুপ্রিন্ট, সংস্করণ ১.০ (২০২৬)। CC BY 4.0।}
\end{minipage}
\end{document}
